% =========================================================
\chapter{Lecture 7 : Authentication}
% =========================================================

\begin{rembox}[title=In a nutshell]
\textbf{What is this lecture about?} Access control (Lectures 3 and 4) decides what a
\emph{principal} may do, but it assumed we already knew \emph{who} the principal is.
\textbf{Authentication} is that missing first step: verifying a claimed identity. We see the three
classic ways to prove who you are: \textbf{what you know} (passwords: how to transfer, store and
check them safely), \textbf{what you are} (biometrics: enrollment, verification, false positives
and negatives) and \textbf{what you have} (tokens, two-factor authentication), and finally how
\textbf{machines} authenticate with keys and signatures. \\
\textbf{Why is it useful?} If authentication fails, every permission of the victim goes to the
attacker: the best access control in the world protects nothing. Password databases leak every
year, replay and timing attacks are classic exam material, and ``don't design your own
authentication protocol'' is a rule you will need in practice.
\end{rembox}

\section{What is authentication?}

\begin{defnbox}[title=Authentication]
\textbf{Authentication} is the process of \textbf{verifying a claimed identity}: the process
that entities use to verify that other entities are who they claim to be. (Rick says ``Listen
Morty, I am the real Rick''; Morty wonders: is this the real Rick?)
\end{defnbox}

\begin{warnbox}[title=Authentication $\neq$ message authentication]
Do not confuse it with \textbf{message authentication} (Lecture 6: MACs, signatures), which
verifies that a \emph{message}
\begin{itemize}
  \item comes from the designated sender (\textbf{provenance}), and
  \item has not been modified (\textbf{integrity}).
\end{itemize}
Here we authenticate an \textbf{entity} (a user, a machine), not a message. Message
authentication is a \emph{tool} that authentication protocols use.
\end{warnbox}

\begin{defnbox}[title=Where does authentication fit?]
A user asks ``write on file \texttt{foo.txt}?''. Two steps happen before the answer ``Write!'':
\begin{enumerate}
  \item \textbf{Authentication}: the system \textbf{binds the actor to a principal} (the abstract
        entity that is authorized to act: users, connections, processes).
  \item \textbf{Authorization} (access control, Lectures 3 and 4): the system decides whether
        \emph{that principal} is authorized.
\end{enumerate}
The lecture mostly covers the authentication of \textbf{users}, and a bit of \textbf{machine}
authentication.
\end{defnbox}

\begin{defnbox}[title=Ways to prove who you are]
\textbf{Traditional} (the three factors):
\begin{itemize}
  \item \textbf{What you know} --- a secret only you know (password, secret key): proving
        knowledge of the secret proves it is you.
  \item \textbf{What you are} --- a trait inherent to you and to no one else (biometrics).
  \item \textbf{What you have} --- an object only you own (smart card, secure token): proving
        possession proves it is you.
\end{itemize}
\textbf{Modern} (traits that are often unique, sometimes called \textbf{soft biometrics}):
\begin{itemize}
  \item \textbf{Where you are} --- location, IP address: apps learn your usual places and raise an
        alarm when you are far away;
  \item \textbf{How you act} --- behavioural authentication: how you type, move the mouse, press
        the keys;
  \item \textbf{Who you know} --- social ties: your network of friends is nearly unique;
  \item many others. Banking apps already use all of them, and more.
\end{itemize}
\end{defnbox}

\section{What you know: passwords}

\begin{defnbox}[title=Password]
A \textbf{secret shared between the user and the system}. The user has a secret password; the
system checks it to authenticate the user.
\end{defnbox}
